\section{Elementos de prueba de hipótesis}

En la prueba de hipótesis, asumimos una premisa inicial (generalmente relacionada con el valor del estimador) denominada \emph{hipótesis nula} y trataremos de ver si es cierta o no aplicando técnicas estadísticas.

Tenemos otra premisa llamada \emph{hipótesis alternativa}, la cuál es la negación de la hipótesis nula.

\subsection{Hipótesis nula vs. alternativa}

Cuando alguien está haciendo una \emph{investigación} cuantitativa para calibrar el valor de un estimador, el \emph{valor conocido} del parámetro se toma como \emph{hipótesis nula}, mientras que el \emph{nuevo valor} encontrado (de la investigación) se toma como la \emph{hipótesis alternativa}.

En nuestro caso (encontrar la edad media de nuestra ciudad), un investigador puede afirmar que la edad es \emph{menor que 35}. Esto puede servir como la \emph{hipótesis nula}.

Si una nueva agencia afirma que es \emph{mayor que 35}, entonces se puede denominar como la \emph{hipótesis alternativa}.

\subsection{Notación formal}

Denotamos las hipótesis como:
\begin{itemize}
	\item $H_0$: Hipótesis nula (afirmación que se asume verdadera inicialmente)
	\item $H_a$ o $H_1$: Hipótesis alternativa (afirmación que contradice a $H_0$)
\end{itemize}

\begin{ejemplo}
	Un fabricante de cereales afirma que el peso promedio de sus cajas es de 500 gramos. Un inspector de calidad sospecha que el peso es menor.

	\begin{align}
		H_0: \mu &= 500 \text{ g} \\
		H_a: \mu &< 500 \text{ g}
	\end{align}
\end{ejemplo}

\subsection{Tipos de pruebas de hipótesis}

Dependiendo de la hipótesis alternativa, las pruebas pueden ser:

\begin{enumerate}
	\item \textbf{Prueba de cola izquierda:}
	\begin{align}
		H_0: \mu &= \mu_0 \\
		H_a: \mu &< \mu_0
	\end{align}

	\item \textbf{Prueba de cola derecha:}
	\begin{align}
		H_0: \mu &= \mu_0 \\
		H_a: \mu &> \mu_0
	\end{align}

	\item \textbf{Prueba de dos colas:}
	\begin{align}
		H_0: \mu &= \mu_0 \\
		H_a: \mu &\neq \mu_0
	\end{align}
\end{enumerate}

\subsection{Errores en pruebas de hipótesis}

Al tomar decisiones basadas en datos muestrales, podemos cometer dos tipos de errores:

\begin{definicion}[Error tipo I]
	Rechazar la hipótesis nula cuando en realidad es verdadera. La probabilidad de cometer este error se denota como $\alpha$ (nivel de significación):
	\begin{align}
		\alpha = P(\text{Rechazar } H_0 | H_0 \text{ es verdadera})
	\end{align}
\end{definicion}

\begin{definicion}[Error tipo II]
	No rechazar la hipótesis nula cuando en realidad es falsa. La probabilidad de cometer este error se denota como $\beta$:
	\begin{align}
		\beta = P(\text{No rechazar } H_0 | H_0 \text{ es falsa})
	\end{align}
\end{definicion}

\begin{observacion}
	La \emph{potencia} de una prueba es la probabilidad de rechazar correctamente una hipótesis nula falsa:
	\begin{align}
		\text{Potencia} = 1 - \beta = P(\text{Rechazar } H_0 | H_0 \text{ es falsa})
	\end{align}
\end{observacion}

\subsection{Potencia y tamaño de muestra}

La \emph{potencia} de una prueba no es un valor fijo: depende del tamaño real del efecto que se desea detectar ($\mu_a - \mu_0$), del tamaño de muestra $n$ y del nivel de significación $\alpha$ elegido. En el diseño experimental, el investigador puede fijar de antemano tanto $\alpha$ como la potencia deseada $1-\beta$ para determinar cuántas observaciones $n$ debe recolectar.

\begin{teorema}[Tamaño de muestra para una potencia objetivo, prueba $Z$ de una cola]
	Sea una prueba de hipótesis sobre la media de una población normal con desviación estándar conocida $\sigma$, de la forma $H_0: \mu = \mu_0$ contra $H_a: \mu > \mu_0$ (o su análoga de cola izquierda). Si se desea que la prueba alcance una potencia $1-\beta$ para detectar un efecto real de tamaño $\mu_a \neq \mu_0$ al nivel de significación $\alpha$, el tamaño de muestra mínimo requerido es:
	\begin{align}
		n = \left( \frac{(Z_\alpha + Z_\beta)\, \sigma}{\mu_a - \mu_0} \right)^2,
		\label{eq:tam_muestra_potencia}
	\end{align}
	donde $Z_\alpha$ y $Z_\beta$ son los cuantiles superiores de la normal estándar correspondientes a $\alpha$ y $\beta$ respectivamente. Para una prueba de dos colas, se reemplaza $Z_\alpha$ por $Z_{\alpha/2}$.
\end{teorema}

\begin{observacion}
	La fórmula \eqref{eq:tam_muestra_potencia} revela relaciones fundamentales de diseño experimental:
	\begin{itemize}
		\item A mayor tamaño de efecto deseado ($\mu_a - \mu_0$ más grande), \textbf{menor} es el tamaño de muestra necesario: los efectos grandes son más fáciles de detectar.
		\item A mayor potencia exigida (por ejemplo, pasar de $1-\beta = 0.80$ a $1-\beta = 0.95$), \textbf{mayor} es el tamaño de muestra requerido, pues $Z_\beta$ crece.
		\item A menor $\alpha$ (una prueba más estricta contra falsos positivos), también \textbf{mayor} es el tamaño de muestra requerido, pues $Z_\alpha$ crece.
	\end{itemize}
	Esto formaliza la intuición de que existe un \emph{trade-off} triangular entre el error Tipo I, el error Tipo II y el costo (tamaño de muestra) de un experimento.
\end{observacion}

